\section{Totient Maximum --- Problem 69}
\label{prob:69}
\noindent\textbf{Source:} \url{https://github.com/KO0I/projectEuler/blob/reorganize/solutions-by-language/solutions/c/totient_maximum_69.c}

\noindent\textbf{Assessment:} Totient-ratio search implemented.

\subsection*{Mathematical goal}
Euler's totient \(\varphi(n)\) counts positive integers at most \(n\) coprime to \(n\). Find \(n\leq10^6\) maximizing \(n/\varphi(n)\). For distinct prime divisors \(p\) of \(n\),
\[
\varphi(n)=n\prod_{p\mid n}\left(1-\frac1p\right).
\]

\subsection*{How the C code works}
The helper \verb|phi| begins at \verb|res=i| and, whenever a prime factor is found, updates \verb|res -= res/p|, removing repeated copies of that factor from the working integer. \verb|main| loops over candidate \verb|n| values from 2 up to one million, calculates the floating-point ratio, and records the best position in \verb|maxpos|.

\subsection*{Correctness and limitations}
The factorization-based totient computation is appropriate, and the expected maximizing integer is 510510. The code recomputes \verb|phi(n)| after deciding the ratio improved, doubling work for those iterations. Its diagnostic print compares two equal values because \verb|maxval| was just updated. A product of successive small primes offers a faster mathematical route. The printed maximum remains valid if the helper is evaluated correctly.


\subsection*{Correction to historical C}
{\color{red}\textbf{Correction:} The corrected implementation multiplies successive primes while staying at most one million. The original totient factorization was useful but recomputed values unnecessarily.}
